\section[Rédaction Bourbaki n\textsuperscript{o}\,128 : espaces bornologiques et tonnelés]{Rédaction Bourbaki n\textsuperscript{o}\,128 : espaces bornologiques et espaces tonnelés\protect\verifie{Redaction128}}\label{sec:bourbaki-128}

La rédaction n\textsuperscript{o}\,128 des archives Bourbaki (Archives Henri Poincaré, \href{https://archives-bourbaki.ahp-numerique.fr/items/show/537}{item 537}) est l'état~4 du chapitre~IV des \emph{Espaces vectoriels topologiques}, consacré aux espaces localement convexes métrisables. Une mention d'archive sur la première page du chapitre la donne pour l'exemplaire de Laurent Schwartz, daté de mars 1950. Le chapitre est suivi de compléments sur deux théorèmes, l'un de Grothendieck, sur la complétion du dual d'un espace localement convexe, l'autre de Banach, puis de \q{Commentaires}. Ces commentaires sont écrits à la première personne et ne sont pas signés ; rien sur l'exemplaire ne nomme leur auteur, et on ne les attribue ici à personne. Sa transcription reste locale pour l'instant ; on n'en cite ici aucun passage.

\emph{Borné et continu.} Pour une application linéaire $u$ entre espaces normés, être continue et envoyer la boule unité sur une partie bornée sont une seule et même condition. Dans un espace localement convexe quelconque, une partie $B$ est \emph{bornée} si tout voisinage $V$ de $0$ l'absorbe, c'est-à-dire si $B \subset \lambda V$ pour $\lambda$ assez grand. Une application linéaire continue envoie toujours les bornés sur des bornés, mais la réciproque peut échouer. Sur l'espace de Hilbert $\ell^2$, les parties faiblement bornées sont bornées en norme (Banach--Steinhaus), de sorte que l'identité de $\ell^2$ muni de la topologie faible dans $\ell^2$ muni de la norme envoie les bornés sur des bornés ; elle n'est pourtant pas continue. La topologie faible a trop peu de voisinages pour ses bornés.

Les Commentaires isolent la propriété qui rétablit l'équivalence. Un espace localement convexe y est dit \q{bornographique} quand toute partie convexe qui absorbe tous les bornés est un voisinage de~$0$ ; on dit aujourd'hui \emph{bornologique}. La proposition~A affirme qu'un espace localement convexe métrisable est bornologique. Sa démonstration est un argument de suite : si $V$ absorbe les bornés sans être un voisinage de~$0$, on prend dans le $n$-ième voisinage $V_n$ d'une base décroissante un point $x_n$ de $V_n/n$ hors de $V$. La suite $(nx_n)$ tend vers~$0$, donc est bornée, et $V$ devrait l'absorber : pour $\lambda$ fixé et $n \geq \lambda$, on aurait $x_n \in (\lambda/n)V \subset V$, puisque $V$ est convexe et contient~$0$. La proposition~A$'$ en tire la conséquence : sur un espace bornologique, une application linéaire qui envoie les bornés sur des bornés est continue, car l'image réciproque d'un voisinage convexe de~$0$ est une partie convexe qui absorbe les bornés.

\emph{Les tonneaux.} Un \emph{tonneau} est une partie convexe, cerclée, fermée et absorbante. Dans un espace normé, la boule unité fermée en est un. Les Commentaires appellent \q{fort} un espace localement convexe où tout tonneau est un voisinage de~$0$ ; on dit aujourd'hui \emph{tonnelé}. La proposition~B affirme qu'un espace localement convexe de Baire est tonnelé. Si $V$ est un tonneau, il absorbe chaque point, donc $E$ est la réunion des fermés $nV$ ; par le théorème de Baire, l'un d'eux, et donc $V$, a un point intérieur $x$. Comme $V$ est cerclé, $-x$ est aussi intérieur, et par convexité $0 = \tfrac12 x + \tfrac12(-x)$ l'est. Les espaces tonnelés sont le cadre naturel du théorème de Banach--Steinhaus, objet de la proposition~B$'$ : une famille d'applications linéaires continues simplement bornée y est équicontinue.

\begin{encadre}
\begin{enonce}[Bourbaki, rédaction n\textsuperscript{o}\,128, Commentaires, propositions~A, A$'$ et~B, dans la forme de mathlib]
Soient $E$ et $F$ des espaces vectoriels topologiques, sur un corps valué non discret $\mathbb K$ pour $E$ et un corps valué $\mathbb K'$ pour $F$, et $\sigma : \mathbb K \to \mathbb K'$ un morphisme isométrique.
\begin{enumerate}
\item[a)] Si $E$ admet en $0$ une base dénombrable de voisinages, en particulier s'il est métrisable, toute application $\sigma$-semi-linéaire $E \to F$ qui envoie les parties bornées sur des parties bornées est continue.
\item[b)] Si $E$ est un espace de Baire, il est tonnelé au sens de mathlib : toute semi-norme semi-continue inférieurement sur $E$ est continue.
\end{enumerate}
\end{enonce}
\end{encadre}

\noindent\emph{Vérification en Lean.} Les deux énoncés sont déjà dans mathlib, et le fichier ne fait que les rapporter à la page. Le premier est \texttt{LinearMap.\allowbreak{}continuous\_of\_locally\_bounded} (\texttt{propA'}). Mathlib n'a pas la notion d'espace bornologique : les propositions~A et~A$'$ y sont réunies en un seul énoncé, sans la propriété intermédiaire, et sans hypothèse de convexité locale. Le second est l'instance \texttt{BaireSpace.\allowbreak{}instBarrelledSpace} (\texttt{propB}). Mathlib définit les espaces tonnelés par les semi-normes. Pour un espace localement convexe réel ou complexe, cette définition équivaut à celle des Commentaires, car les tonneaux sont exactement les boules unités fermées des semi-normes semi-continues inférieurement ; cette équivalence n'est pas démontrée dans le fichier. Les noms sont ceux de \href{\leanblobmain/Redaction128.lean}{\texttt{Redaction128.\allowbreak{}lean}}.

\begin{remarque}
Les Commentaires notent que les espaces (LF), limites inductives strictes de suites d'espaces de Fréchet, sont bornologiques et tonnelés sans être métrisables ni de Baire. Les deux notions dépassent donc les cas que traitent les propositions~A et~B.
\end{remarque}

Ne sont pas démontrés ici : les remarques sur les espaces (LF) ; la proposition~B$'$, que mathlib démontre pour les espaces tonnelés (\texttt{WithSeminorms.\allowbreak{}banach\_steinhaus}) mais que le fichier ne reprend pas ; l'espace $E^f$ des Commentaires et les propositions~C, C$'$ et~C$''$ sur les parties bornées ; les compléments sur les théorèmes de Grothendieck et de Banach qui précèdent les Commentaires.
